\section{Fixed-inertia cubic curves and the actual integral finiteness condition}
\label{sec:cubic-fixed-weight-curve-m2}

This section applies the changed-tuple ordinary chart to the
characteristic-\(269\) curves of the actual fixed-local-type model.
The two deformation problems are kept separate.  The first has
variable ordinary characters; the second has the exact integral
inertia characters proved in AC-G65-flag.  That additional condition
is what makes the weight ideal vanish on the curves considered here.

\paragraph{Actual data and the one retained support hypothesis.}
Put \(p=269\) and \(K=\mathbb Q(\omega)\).  Use the actual cubic
model \(\mathscr S\) over its original coefficient DVR
\(\mathcal O_0\), its universal semidirect representation
\(r_{\mathscr S}\), and its weighted presentation from AC-G65-model
and AC-G66.  Retain their original datum, coefficient, marked-element,
local-type, and source hypotheses.  Write the residual tuple on
\(G_K\) as
\[
 \overline{\boldsymbol a}
   =(\bar\chi^{\,j}\bar\gamma,1,
                  \bar\chi^{\,j}\bar\gamma^{-1}),\qquad j=1,2.
 \tag{G109.1}
\]
Here \(\gamma\) is the original \(\mathcal O_0\)-valued rational
character; bars always mean its finite residual character.
The original marked elements are \(c,g_1\), with
\(r_{\mathscr S}(c).\mathrm{mat}=I\) and nonzero residual marked
difference at \(g_1\).

Let \(H\supset K\) be one fixed actual finite CM host for the
changed-tuple construction AC-G90--AC-G108.  Its prescribed local
maps are the full decomposition groups over \(p\), with their
inertia subgroups.  Write \(\mathcal O\) for the host coefficient
DVR of AC-G94--AC-G108, with residue field \(k\), and write
\(k_0\) for the residue field of \(\mathcal O_0\).
Fix their embeddings into one algebraic closure of \(\mathbb F_p\)
compatibly with the actual residual tuple.  The finite values
in (G109.1)--(G109.2) have their prescribed common identification.
There is no assumed homomorphism between \(\mathcal O_0\) and
\(\mathcal O\).  In particular a ramified native coefficient
choice does not change any original \(e_{\mathcal O_0}=1\)
condition where that condition was imposed.
Let \(N\) be the actual closed normal away-inertia kernel on
\(G_H\), and put \(G=G_H/N\), as in AC-G90-quotient.  Thus the
remaining away-\(p\) inertia has finite order prime to \(p\),
and the character tuple
\[
 \boldsymbol\theta
       =\overline{\boldsymbol a}|_{G_H}
 \tag{G109.2}
\]
descends to \(G\).  Retain the actual separation at every local
map and the host normalization witness \(g_0\), for which
\(\theta_0(g_0)-\theta_2(g_0)\ne0\).

Use \(B\subset A\subset R\) for the host rings of AC-G94--AC-G95,
and \(\kappa_{v,i}\) for their ordered local characters.  In contrast,
write \(\mathscr B\subset\mathscr S\) for the degree-zero ring of
the original model.  The actual condition \texttt{M2} is
\(\mathscr B\) finite over \(\mathcal O_0\).
Define
\[
 \begin{split}
 U&=(r_{ij}(g):i<j)A,\qquad L=(r_{ij}(g):i>j)A,\\
 K_v&=I_v\cap\bigcap_i\ker(\theta_i\circ\operatorname{loc}_v),\\
 \lambda_\theta&=(\varpi)B+
    (\kappa_{v,i}(s)-1:s\in K_v,\ v,i)B.
 \end{split}
 \tag{G109.3}
\]
The support hypothesis used below is the precise conclusion of
AC-G108-propagation: there is an ideal \(I_{\rm off}\subset B\)
such that
\[
 \begin{split}
 \sqrt{I_{\rm off}+\lambda_\theta}&=\mathfrak m_B,\\
 \mathfrak q\in\operatorname{Spec}A,\quad
 U\not\subset\mathfrak q,\quad L\not\subset\mathfrak q
       &\ \Longrightarrow\ I_{\rm off}\subset\mathfrak q\cap B.
 \end{split}
 \tag{G109.4}
\]
Whenever AC-G108 is used to supply (G109.4), all its premises,
including its uniform endpoint support condition \((E)\), remain
premises here.  A single native seed point or a finite-weight
native quotient alone is not a substitute for (G109.4).

\begin{proposition}[AC-G109-curve: actual curves kill the joint-kernel weight ideal]
\label{AC-G109-curve}
Let
\[
 f:\mathscr S\longrightarrow V
 \tag{G109.5}
\]
be a local \(\mathcal O_0\)-algebra homomorphism to a complete
discrete valuation ring with finite residue field, with
\(\varpi_0 V=0\).  There is an injective extension
\(V\hookrightarrow V'\) of equal-characteristic complete DVRs,
a complete Noetherian local \(\mathcal O\)-subalgebra
\(C\subset V'\) with residue field exactly \(k\), and a
strict constant \(J\)-orthogonal change of basis, such that
the restricted curve on \(G_H\) descends to a
normalized ordinary lift over \(C\) in the exact functor of
AC-G94.  Its universal specialization
\[
             R\longrightarrow C\subset V'
 \tag{G109.6}
\]
annihilates \(\lambda_\theta\).  No irreducibility assumption on
the curve or on its restriction is required.
\end{proposition}

\begin{proof}
\emph{1. The test algebra and the away kernel.}
The residue field of \(V\) is a finite extension \(k'/k_0\).
Simple-root Hensel lifting gives its canonical constant field
inside the equal-characteristic complete DVR.  Choose a
uniformizer \(t\); then \(V=k'[[t]]\).
Choose a finite field \(k''\) in the fixed algebraic closure
containing both \(k'\) and \(k\).  Coefficientwise inclusion gives
\(V\hookrightarrow V'=k''[[t]]\).  Set
\[
                 C=k+tV'\subset V'.
 \tag{G109.7}
\]
This is a local ring with maximal ideal \(tV'\), residue field
\(k\), and \((tV')^n=t^nV'\).  It is finite as a module over
\(k[[t]]\): a basis of \(k''/k\), with the nonconstant basis
vectors multiplied by \(t\), supplies finite generators.
Thus it is Noetherian and complete, with exactly the induced
topology from \(V'\).  Its \(\mathcal O\)-algebra structure is
the prescribed reduction \(\mathcal O\to k\).

After this injection every entry of the restricted representation
has its prescribed residual value in \(k\), so lies in \(C\);
its inverse matrices do also.  Continuity in \(V'\) is continuity
in \(C\).  Only this representation is placed over \(C\):
no homomorphism \(\mathscr S\to C\) over \(\mathcal O_0\),
no map \(\mathcal O_0\to\mathcal O\), and no base-change theorem
for the original deformation ring are needed.
At an away-\(p\) inertia
element in the joint residual kernel the original modified
\texttt{Kill} condition gives the identity matrix.  Restriction
to \(H\) retains this assertion.  Continuity then kills the
closed normal subgroup generated by all such elements, exactly
\(N\).  This proves descent to \(G\), without inflating an
\(H^2\)-class or replacing this quotient by an unrestricted
Galois group.

The original equality at \(c\) gives the exact polarization
\[
       r(g^\star)=Jr(g)^tJ,\qquad g^\star=cg^{-1}c^{-1}.
 \tag{G109.8}
\]
It descends because the actual kernel is stable under \(\star\).

\emph{2. Integral normalization at the host witness.}
The old marked element need not equal the host witness.
Put \(Q=r(g_0)-r(g_0^\star)\in M_3(C)\).
Equation (G109.8) gives \(Q^tJ=-JQ\).  Since \(2\) is a unit,
its trace and determinant are zero and its characteristic
polynomial is
\[
       X(X^2-b),\qquad b=\tfrac12\operatorname{Tr}(Q^2).
 \tag{G109.9}
\]
Its reduction is the diagonal matrix
\(\operatorname{diag}(a_0,0,-a_0)\), where
\(a_0=\theta_0(g_0)-\theta_2(g_0)\ne0\).
Hensel lifting gives a unique \(a\in C^\times\) reducing to
\(a_0\) with \(a^2=b\).  The Lagrange projectors for the
three roots \(a,0,-a\) are defined over \(C\).
Their coordinate columns form a matrix \(P\equiv I\)
diagonalizing \(Q\), exactly as in AC-G94.

The \(J\)-skew identity implies that eigenlines whose
eigenvalues have unit sum are orthogonal.  Consequently
\[
            P^tJP=
       \begin{pmatrix}0&0&u\\0&v&0\\u&0&0\end{pmatrix},
            \qquad u,v\in1+\mathfrak m_C.
 \tag{G109.10}
\]
Multiply the first column by \(u^{-1}\) and the middle
column by the unique square root \(v^{-1/2}\in1+\mathfrak m_C\).
The resulting matrix \(h\equiv I\) satisfies
\[
 h^tJh=J,\qquad h^{-1}Qh=\operatorname{diag}(a,0,-a).
 \tag{G109.11}
\]
Replacing \(r\) by \(h^{-1}rh\) preserves (G109.8), the
residual tuple, continuity and all invariant flag conditions.
It gives the literal off-diagonal normalization at \(g_0\).
This argument uses only separation and polarization, and
therefore applies equally to reducible curves.

\emph{3. The full local flags and their fixed inertia.}
AC-G65-flag holds over every local test algebra, explicitly
including the original characteristic-\(269\) ring \(V\) here.  It gives
an upper flag on \(G_{K_p}\) with ordered inertia characters
\(\chi^j\gamma,1,\chi^j\gamma^{-1}\).
After reduction to \(V\) and injection into \(V'\), these are
the constant characters in (G109.1).  Restrict the flag to each
actual full host local
group.  The compatible local embeddings give the residual
labels (G109.2).  For a conjugate embedding, transporting by
the global representation preserves these labels and makes
the transported basis strict after multiplication by its
residual diagonal.  The conjugate-dual flag gives the same
assertion at the paired place.

Equivalently, one may form the labeled spectral projectors
of a local residual-separating element directly over \(C\).
They belong to \(M_3(C)\), their coordinate matrix is strict,
and over \(V'\) this is the same ordered flag by uniqueness.
Thus the local flag and all its quotient characters are
defined over \(C\), even if a first flag basis was chosen
only over \(V'\).  Their values reduce to \(k\), are continuous,
and satisfy the exact ordered-product equations on the full
local group.  The strict global change (G109.11) leaves those
quotient characters unchanged.

Generic framed universality and the exact equations in
AC-G94 now give (G109.6).  Its local characters on inertia
are precisely \(\theta_i\circ\operatorname{loc}_v\), by the
uniqueness part of that theorem.  Every generator
\(\kappa_{v,i}(s)-1\), \(s\in K_v\), therefore vanishes.
The coefficient uniformizer vanishes as well.  This proves
the asserted annihilation of \(\lambda_\theta\).
\end{proof}

\begin{theorem}[AC-G109-M2: off-triangular support supplies actual integral finiteness]
\label{AC-G109-M2}
Under the actual setup above and (G109.4), every curve
(G109.5) has the full original cubic seed characteristic
polynomial on \(G_K\).  In the original normalized basis
it is upper or lower triangular and has the prescribed
seed diagonal.  Consequently
\[
                 \mathscr B\text{ is finite over }\mathcal O_0;
                 \qquad\texttt{M2 holds for the actual model.}
 \tag{G109.12}
\]
In particular there is no further prescribed-curve Hecke
containment premise between the actual conclusion of
AC-G108 and this \texttt{M2} conclusion.  This is a conditional
implication from (G109.4), not a removal of the uniform
endpoint support premise needed to establish that hypothesis.
\end{theorem}

\begin{proof}
\emph{1. A balanced monomial forces triangularity on the host.}
Fix a curve and use AC-G109-curve.  Let
\(\xi:A\to C\subset V'\) be its specialization and
\(\mathfrak q=\ker\xi\), a prime.  Suppose it contains
neither \(U\) nor \(L\).
Then (G109.4), together with \(\lambda_\theta\subset\ker\xi\),
gives \(\mathfrak m_B\subset\ker(\xi|_B)\).
There is a surviving positive-weight matrix entry \(x\)
of weight \(a>0\), and a surviving negative-weight entry
\(y\) of weight \(-b<0\).  The element
\[
                        x^b y^a\in A_0=B
 \tag{G109.13}
\]
has zero residue and hence lies in \(\mathfrak m_B\).
Its image is nonzero in the domain \(V'\), a contradiction.
Therefore the host representation is upper or lower triangular
in the normalized basis of AC-G109-curve.

\emph{2. Identify the diagonal on every coefficient-prime inertia.}
Write \(a_i:G\to C^\times\) for its diagonal characters.
They reduce to \(\theta_i\).
In the triangular basis each local residual-separating matrix
is triangular.  Its labeled Hensel projectors are polynomial
in that matrix, so their coordinate-column matrix is
triangular of the same kind, invertible and congruent to one.
Conjugation by an invertible upper triangular matrix preserves
the diagonal of an upper triangular matrix; the same holds
for lower triangular matrices.  In this canonical local basis
AC-G94 identifies the diagonal with the ordered local
characters.  Consequently
\[
       a_i|_{I_v}=\theta_i|_{I_v}\quad\text{for all }v\mid p.
 \tag{G109.14}
\]
This reasoning also covers the lower-triangular case: its
canonical local matrix is lower triangular, while the ordinary
equations give the upper flag, and its diagonal is unchanged.
It does not identify characters merely by an unordered
characteristic polynomial.

\emph{3. The remaining principal-unit characters are trivial.}
Regard the finite characters \(\theta_i\) as constant
characters with values in \(k^\times\subset C^\times\), and put
\(\nu_i=a_i\theta_i^{-1}\).
These are continuous characters with values in
\(1+\mathfrak m_C\), trivial on all inertia over \(p\)
by (G109.14).
Away from \(p\), the actual quotient has inertia factoring
through a finite group of order prime to \(p\).
Its image under \(\nu_i\), reduced modulo any
\(\mathfrak m_C^n\), lies in a finite \(p\)-group and is
therefore trivial.  Separation of \(C\) makes \(\nu_i\)
trivial on that inertia as well.

Thus every finite reduction of \(\nu_i\), viewed on \(G_H\),
is an everywhere unramified abelian \(p\)-character.
The actual class-exponent argument of AC-G99-frame applies:
these characters factor through the finite ordinary ideal
class group of \(H\).  There are no real places, since \(H\)
is CM.  A single positive integer \(h\), for example the
exponent of its \(p\)-class group, kills every finite reduction.
Passing to all \(n\) gives \(\nu_i^h=1\).

The principal-unit group of a characteristic-\(p\) domain
has no torsion.  Indeed write \(h=p^s m\), with \(p\nmid m\).
If \(u\in1+\mathfrak m_C\) and \(u^h=1\), then
\((u^m-1)^{p^s}=0\), so \(u^m=1\).
The second factor in
\((u-1)(1+u+\cdots+u^{m-1})=0\)
is a unit, having nonzero residue \(m\).  Hence \(u=1\).
It follows that \(a_i=\theta_i\) for every \(i\).
We have proved
\[
 \det(X-r_f(g).\mathrm{mat})
       =\prod_{i=0}^2(X-\overline a_i(g))
                  \quad(g\in G_H).
 \tag{G109.15}
\]
Characteristic polynomials are unchanged by the strict
host normalization.  Thus (G109.15) holds for the image of
the original curve in \(V'\), and descends to \(V\) by
injectivity, since its prescribed coefficients already lie
in the original constant field \(k_0\).  No upper bound for
the class-group exponent was used.

\emph{4. Descend the constant polynomials through the finite host.}
For each \(g\in G_K\), finite index gives an integer \(n>0\)
such that \(g^n\in G_H\).  The host need not be normal for
this elementary assertion.  Since the multiplier is one on
\(G_K\), \(r_f(g^n).\mathrm{mat}=r_f(g).\mathrm{mat}^{\,n}\).
By (G109.15) every eigenvalue of this power is in \(k_0^\times\).
Each eigenvalue of \(r_f(g).\mathrm{mat}\) is therefore a
root of \(X^n-a\) for some \(a\in k_0^\times\), and is algebraic
over \(k_0\).  This argument permits \(p\mid n\).
The characteristic coefficients are algebraic over \(k_0\),
belong to \(V\), and reduce to the prescribed coefficients
of the original seed in \(k_0\).

An element \(z\in V\) algebraic over \(k_0\), with residue
\(z_0\in k_0\), equals \(z_0\): the ring \(k_0[z]\) is a finite
field, and if \(z-z_0\ne0\) its inverse in that field would
belong to \(V\), contradicting \(z-z_0\in\mathfrak m_V\).
Thus the full seed polynomial identity holds on all of \(G_K\).
This is finite-index descent of characteristic coefficients,
and does not require the restricted representation to be
irreducible or semisimple.

\emph{5. Original normalization gives the required coordinate vanishing.}
We make explicit the extension of AC-G59-polynomials to this
coefficient prime.  Its proof uses only characteristic different
from two, the original normalization at \(c,g_1\), and nonzero
marked residual difference.  All are retained by the actual
AC-G65 model.  The literal generic theorem
\path{UnitaryThreeGroups/ExistsCoordOfNotMatIrr.lean},
\texttt{shape\_cases}, therefore gives Type I, upper, lower,
or Type II over an algebraic closure of \(\operatorname{Frac}V\).

An irreducible representation of dimension greater than one
cannot have the same characteristic polynomials as a sum of
characters.  For completeness, the image algebra on the direct
sum of the irreducible representation and those characters is
finite dimensional and semisimple: it acts faithfully on a
semisimple module, so its radical is zero.  Its matrix factor
for the irreducible representation is distinct from the
one-dimensional factors.  An element acting there as \(E_{11}\)
and as zero on the character factors is a linear combination
of group operators.  The assumed trace identity would give
\(1=0\).  This excludes Type I.
In Type II the middle character is a seed-polynomial root
with residue one, so the preceding finite-field argument makes
it one.  Canceling \(X-1\) gives the irreducible endpoint block
the two endpoint character polynomials, the same contradiction.
Hence the original curve is upper or lower triangular.
Vanishing of its entries over the algebraic closure is
vanishing already in \(V\).  Each diagonal entry is now a
root of the seed polynomial, is algebraic over \(k_0\), and
has its prescribed residual label.  It equals that label.

For the original weighted presentation, all coordinates are
embedded universal matrix-entry differences, as in AC-G51
and its exact AC-G65-model extension.  Consequently every
weight-zero coordinate maps to zero, and one entire sign
of the nonzero-weight coordinates maps to zero.  Matrices
in the other multiplier coset give the same entries by right
multiplication by \(c\), whose matrix is one.

\emph{6. Apply the integral characteristic-\(p\) curve criterion.}
The model carries the complete homogeneous quotient
presentation used in AC-G56 and AC-G66.  Its coordinates
belong to the maximal ideal; its formal weight-zero part
is precisely \(\mathscr B\).  The generic source theorem
\path{CurveDVRSpecialization/FinCrit.lean},
\texttt{UpperTheory.V.Pr.finCrit}, requires a complete
coefficient DVR with finite residue field, an \texttt{IsCO}
target, all presentation coordinates in its maximal ideal,
and exactly the preceding two coordinate-vanishing assertions
for every \texttt{IsCD} characteristic-\(p\) curve.
All these hypotheses have now been verified for \(\mathscr S\).
That theorem is applied over the original \(\mathcal O_0\)
and gives finiteness of the formal degree-zero part over
\(\mathcal O_0\), which is (G109.12).
Its integral proof detects nonnilpotent zero-weight monomials
modulo \(\varpi_0\) on complete-DVR curves and then uses complete
module Nakayama; thus vertical components are included.
No generic-fiber finiteness inference is involved.
\end{proof}

\paragraph{Exact source interface and scope.}
The complete source proofs \path{FinPtOf.lean},
\path{FinEndOfFinPt.lean}, \path{GfResZpow.lean},
\path{FinCrit.lean}, and
\path{ModelIdealFinitenessOfWeightZeroFinite.lean} were read
for this interface.  The old \texttt{finEnd\_of\_finPt} explicitly
uses a single-character \texttt{GalLink} and the identity with
\(\Psi^{w_i}\); it is not applied unchanged here.
AC-G109-curve supplies the actual changed-tuple map instead.
The balanced monomial and class-exponent arguments above are
the direct replacement of the old \texttt{FinPt} implication.
Only the genuinely generic shape and integral \texttt{finCrit}
statements are then applied.

This closes the passage from actual off-triangular support to
\texttt{M2} on the fixed-local-type cubic model.  It supplies no
uniform endpoint support theorem, no rank-four singular regulator,
and no arbitrary varying-weight curve containment.  In particular,
\(\kappa_{v,i}(s)-1\) need not vanish for \(s\in K_v\) on a
general ordinary deformation.  It vanishes here because
AC-G65-flag fixes the actual inertia characters before the
characteristic-\(269\) curve is chosen.

\paragraph{Geometric dictionary.}
The finite host is a finite-cover restriction of the same residual
local system.  The strict orthogonal basis change selects the
normalized ordinary chart without changing its local flag
characters.  The balanced monomial argument rules out a curve
with both upper and lower directions after its invariant weight
fiber has become the closed point.  The final complete-DVR test
is an integral special-fiber finiteness statement.  This uses the
previously audited finite-cover dictionary and the explicit
algebra above; no new archival Grothendieck attribution is made.
